%% main.tex -- CARL (MLSys 2027 submission)
%%
%% Format: MLSys 2-column (same submission format as MLSys 2026).
%% Requires the MLSys style file `mlsys2024.sty` (and `fancyhdr`) to be present
%% in this directory or on TEXINPUTS. The MLSys author kit reuses the same
%% style year-to-year; drop the current year's .sty in here to compile.

\documentclass{article}

% Recommended, but optional, packages for figures and better typesetting:
\usepackage{microtype}
\usepackage{graphicx}
\usepackage{booktabs}   % for professional tables
\usepackage{amsmath}
\usepackage{amssymb}
\usepackage{url}

% hyperref makes hyperlinks in the resulting PDF.
\usepackage{hyperref}

% Attempt to make hyperref and algorithmic work together better:
\newcommand{\theHalgorithm}{\arabic{algorithm}}

% Use the following line for the initial blind version submitted for review:
\usepackage{mlsys2024}

% If accepted, instead use the following line for the camera-ready submission:
%\usepackage[accepted]{mlsys2024}

% The \mlsystitle you define below is probably too long as a header.
% Therefore, a short form for the running title is supplied here:
\mlsystitlerunning{Adaptive Online Scheduling for Non-Stationary LLM Serving}

\begin{document}

\twocolumn[
\mlsystitle{CARL: Adaptive Online Scheduling, Not Speculative\\Decoding, Drives LLM Serving Under Workload Change}

% It is OKAY to include author information, even for blind
% submissions: the style file will automatically remove it for you
% unless you've provided the [accepted] option to the mlsys2024
% package.
\begin{mlsysauthorlist}
\mlsysauthor{Anonymous Author(s)}{anon}
\end{mlsysauthorlist}

\mlsysaffiliation{anon}{Anonymous Institution}

\mlsyscorrespondingauthor{Anonymous Author(s)}{anon@example.com}

% You may provide any keywords that you find helpful for describing your paper.
\mlsyskeywords{LLM serving, online scheduling, contextual bandits, non-stationary workloads, speculative decoding}

\vskip 0.3in

\begin{abstract}
Modern LLM serving systems are configured offline: an operator
grid-searches batch size, chunk size, speculative-decoding depth, and
eviction policy on a representative trace and freezes the winner. We
show this is the wrong abstraction for production traffic, which is
non-stationary --- alternating between latency-sensitive interactive
bursts and throughput-oriented batch phases whose optimal
configurations are nearly disjoint (interactive: batch=4;
batch: batch=32 with preemption disabled). On such workloads a per-model-tuned static
configuration sacrifices 9--10\% throughput and inflates TTFT p99 by
$\sim$8$\times$ (18,058~ms vs 2,199~ms on TinyLlama), and a naive
reactive autotuner does worse than static (18.9\% vs 16.2\% oracle gap)
by thrashing without regime structure. We present CARL, a
contextual-bandit scheduler that re-selects the serving operating point
online as the regime shifts. CARL reaches 108.2~$\pm$~3.4\% of a
per-model-tuned static best across two model families with zero
retuning (real-GPU, +10.6\% TinyLlama, +5.8\% Qwen2), and its LinUCB
policy closes the per-regime oracle gap to 0.45\% versus 16.2\% for
static. Crucially, we isolate where this gain comes from: across every
stationary single-regime workload CARL is statistically identical to
the static best (Cohen's $d = 0$, $p = 1.0$), and the entire advantage
appears only under regime change (NON-STATIONARY $d = 23.2$,
$p = 2.3\times10^{-41}$, $n=30$). Because the serving engine is held
fixed across both conditions, we conclude that adaptive online
scheduling, rather than speculative decoding, is the primary driver of
performance improvements under changing LLM-serving workloads. We
further delimit the claim: on stable, sequential, or memory-bound
traffic, online exploration is a net cost (CARL loses 5/5 stress
workloads, up to $-$24.2\%), making regime non-stationarity the precise
precondition for adaptive serving.
\end{abstract}
]

% This command actually creates the footnote in the first column
% listing the affiliations and the copyright notice. The command takes
% one argument, which is text to display at the start of the footnote.
\printAffiliationsAndNotice{}

\section{Introduction}
% Argues: production LLM traffic is non-stationary, offline-frozen static configs lose 9-10% throughput and ~8x TTFT p99 under regime shift, and the central thesis that adaptive online scheduling (not speculative decoding) is the driver; states contributions.

\section{Background and Motivation}
% Argues: explains the serving config space (batch size, chunk size, speculative decoding, eviction), how static-best and reactive AutoTuner baselines work, and why neither tracks interactive<->batch regime shifts; introduces contextual bandits.

\section{CARL Design}
% Argues: describes CARL's contextual-bandit scheduler --- context features, the arm/config space, the LinUCB policy and regime detection, and how it re-selects the operating point online without modifying the serving engine.

\section{Evaluation}
% Argues: experimental setup and headline results --- real-GPU cross-model gains (108.2% of static-best, +10.6% TinyLlama, +5.8% Qwen2), throughput/TTFT wins on non-stationary traffic, and statistical validation (n=30 paired t-test, large effect sizes).

\subsection{Design Scope and Limitations}
\label{sec:limitations}

Speculative decoding in our implementation is single-request only: the
scheduler runs the greedy self-speculative (early-exit) path solely when
exactly one request is in the decode phase and reverts to standard
batched decoding whenever two or more requests decode concurrently, and
the implemented draft/target decoder is not wired into the paged-cache
scheduler --- so speculation is inactive across the batched regime and is
held disabled in the experiments reported here.

\section{Where the Gain Comes From}
% Argues: isolates the source of improvement --- zero CARL-vs-static delta in every stationary regime (d=0, p=1.0), the gain appearing only under regime change, near-oracle LinUCB gap (0.45%), and the failure-case boundary (5/5 losses on stable/sequential/memory-bound stress workloads).

\section{Related Work}
% Argues: positions CARL against prior LLM serving systems (continuous batching, speculative decoding, paged attention), offline/auto config tuning, and bandit/RL approaches for systems control.

\section{Conclusion}
% Argues: restates that adaptive online scheduling is the primary driver under non-stationary workloads, that the gain is conditional on regime change, and outlines limitations and future work.


\bibliography{bibliography}
\bibliographystyle{mlsys2024}

\end{document}
